\phantomsection
\addcontentsline{toc}{subsection}{Inciso b)}
\subsection*{Inciso b)}

% Esta seccion cita muchos nombres de clase largos, que no se pueden partir con
% guion: se prefiere estirar los espacios antes que desbordar el margen.
\sloppy

\textit{Implementa con 4 tipos de elementos (CuentaAhorro, CuentaCorriente, CuentaInversión,
CuentaEmpresarial) y 3 operaciones (EjecutarTransacción, RevertirTransacción,
AuditarTransacción). Cada operación debe comportarse de manera diferente según el tipo de
cuenta.}

\vspace{0.4cm}

\subsubsection*{Qué se implementó}

La implementación completa, compilable y ejecutable, está en
\texttt{docs/CP4/codigo/ejercicio12/} (Java 17, sin dependencias externas). Comprende las
cuatro cuentas, las tres operaciones como visitantes concretos y la infraestructura que exige
la restricción de concurrencia:

\begin{itemize}[leftmargin=1.2cm]
    \item \textbf{Elementos:} \texttt{ElementoContable} (interfaz con \texttt{accept}),
          \texttt{Cuenta} (clase base con el saldo, el cerrojo propio y el historial) y las
          cuatro cuentas concretas.
    \item \textbf{Operaciones:} \texttt{EjecutarTransaccionVisitor},
          \texttt{RevertirTransaccionVisitor} y \texttt{AuditarTransaccionVisitor}, más
          \texttt{CalcularInteresesVisitor} añadido al final como comprobación empírica de la
          extensibilidad.
    \item \textbf{Infraestructura:} \texttt{VisitanteTransaccional} (decorador que toma el
          cerrojo de la cuenta visitada), \texttt{Transferencia} (orden total de adquisición
          de dos cerrojos) y \texttt{Demo} (las ocho comprobaciones cuya salida se reproduce
          más abajo).
    \item \textbf{Valores inmutables:} \texttt{Transaccion}, \texttt{Movimiento},
          \texttt{ResultadoOperacion} e \texttt{InformeAuditoria} como \textit{records}.
\end{itemize}

El dinero se representa en centavos con \texttt{long}, nunca con \texttt{double}, para que las
comprobaciones de los saldos sean exactas.

\subsubsection*{Las doce combinaciones}

El enunciado pide que cada operación se comporte de manera diferente según el tipo de cuenta.
Las reglas implementadas son éstas, para las operaciones EjecutarTransacción,
RevertirTransacción y AuditarTransacción. Ninguna es cosmética: cada celda cambia el
resultado.

\begin{table}[h!]
    \centering
    \small
    \begin{tabular}{p{2.1cm}p{3.55cm}p{3.55cm}p{3.55cm}}
        \toprule
        \textbf{Cuenta} & \textbf{Ejecutar} & \textbf{Revertir} & \textbf{Auditar} \\
        \midrule
        Ahorro & Rechaza cualquier salida que deje el saldo en negativo; cobra 1,50 a partir
                 del cuarto retiro del mes.
               & Devuelve el principal, reembolsa la comisión y descuenta el retiro contado.
               & Riesgo MEDIO y hallazgo si se excedió el límite de retiros del mes. \\
        \addlinespace
        Corriente & Admite sobregiro hasta la línea de crédito y cobra 5,00 al usarla.
                  & Devuelve el principal y reembolsa la comisión de sobregiro.
                  & Riesgo MEDIO o ALTO según la profundidad del sobregiro frente a la línea. \\
        \addlinespace
        Inversión & Toda salida liquida participaciones y paga el 2\% de penalización.
                  & Devuelve el principal pero retiene la penalización: no es simétrico.
                  & Riesgo según el peso de las penalizaciones acumuladas sobre el saldo. \\
        \addlinespace
        Empresarial & Exige segunda firma desde 50\,000,00 y marca como reportable todo lo que
                      supere 10\,000,00.
                    & La reversión misma exige la segunda firma; sin ella se rechaza.
                    & Cuenta las operaciones que superan el umbral de reporte regulatorio. \\
        \bottomrule
    \end{tabular}
\end{table}

\subsubsection*{El doble despacho en el código}

{\small
\begin{verbatim}
public interface ElementoContable {
    <R> R accept(VisitanteOperacion<R> v);
}

public final class CuentaAhorro extends Cuenta {
    @Override
    public <R> R accept(VisitanteOperacion<R> v) {
        return v.visitAhorro(this);        // 2do despacho
    }
}

public interface VisitanteOperacion<R> {
    R visitAhorro(CuentaAhorro c);
    R visitCorriente(CuentaCorriente c);
    R visitInversion(CuentaInversion c);
    R visitEmpresarial(CuentaEmpresarial c);
}
\end{verbatim}
}

El parámetro de tipo \texttt{R} no es decorativo: permite que el visitante \emph{devuelva} su
resultado en lugar de acumularlo en un campo, que es la propiedad que lo hace utilizable en
paralelo. \texttt{EjecutarTransaccionVisitor} devuelve \texttt{ResultadoOperacion},
\texttt{AuditarTransaccionVisitor} devuelve \texttt{InformeAuditoria} y
\texttt{CalcularInteresesVisitor} devuelve \texttt{Long}, todos con la misma interfaz de
elemento.

\subsubsection*{Una operación, cuatro comportamientos}

Extracto de \texttt{EjecutarTransaccionVisitor}, donde se ve que la diferencia entre tipos de
cuenta vive en la operación y no en las cuentas:

{\small
\begin{verbatim}
private final Transaccion t;    // el parametro viaja en el visitante

@Override
public ResultadoOperacion visitAhorro(CuentaAhorro c) {
    if (c.saldo() + t.delta() < 0) {
        return ResultadoOperacion.rechazada(c,
                "el ahorro no admite sobregiro");
    }
    long cargos = 0L;
    if (t.tipo() == TipoTransaccion.RETIRO) {
        c.contarRetiro(1);
        if (c.retirosDelMes() > CuentaAhorro.LIMITE_RETIROS_MES) {
            cargos = CuentaAhorro.COMISION_EXCESO;
        }
    }
    c.aplicar(t.delta() - cargos);
    ...
}

@Override
public ResultadoOperacion visitInversion(CuentaInversion c) {
    long cargos = 0L;
    if (t.tipo() != TipoTransaccion.DEPOSITO) {
        cargos = t.montoCentavos()
                * CuentaInversion.PENALIZACION_BPS / 10_000L;
        c.acumularPenalizacion(cargos);
    }
    c.aplicar(t.delta() - cargos);
    ...
}
\end{verbatim}
}

\subsubsection*{Cómo se cumple la restricción}

La restricción se resolvió con los tres mecanismos anunciados en el inciso (a), ahora
implementados:

\textbf{1. Visitantes que devuelven en lugar de acumular.}
\texttt{AuditarTransaccionVisitor} no tiene ni un campo: construye y devuelve un
\texttt{InformeAuditoria} inmutable. Por eso una sola instancia se comparte entre hilos sin
sincronización, como hace \texttt{Demo}. \texttt{EjecutarTransaccionVisitor} sí tiene estado
(la transacción), pero es inmutable y se crea uno por operación.

\textbf{2. El cerrojo vive en la cuenta, no en el visitante.} Cada \texttt{Cuenta} lleva su
\texttt{ReentrantLock} porque la cuenta es la unidad de consistencia del dominio. El decorador
lo toma y lo libera, de modo que ningún visitante concreto sabe que existen cerrojos:

{\small
\begin{verbatim}
public final class VisitanteTransaccional<R>
        implements VisitanteOperacion<R> {

    private final VisitanteOperacion<R> interno;

    @Override
    public R visitAhorro(CuentaAhorro c) {
        return bajoCerrojo(c, () -> interno.visitAhorro(c));
    }

    private R bajoCerrojo(Cuenta c, Supplier<R> accion) {
        c.cerrojo().lock();
        try {
            return accion.get();
        } finally {
            c.cerrojo().unlock();
        }
    }
}
\end{verbatim}
}

Como el cerrojo es por cuenta, las operaciones sobre cuentas distintas no se estorban y las
que afectan a la misma cuenta se serializan: exactamente lo que pide el enunciado.

\textbf{3. Orden total de adquisición en las transferencias.} Una transferencia toca dos
cuentas, luego necesita dos cerrojos y aparece riesgo de interbloqueo. Se toman siempre por
identificador ascendente, nunca en el orden en que llegan los argumentos:

{\small
\begin{verbatim}
Cuenta primera = origen.id().compareTo(destino.id()) <= 0
        ? origen : destino;
Cuenta segunda = primera == origen ? destino : origen;
primera.cerrojo().lock();
try {
    segunda.cerrojo().lock();
    ...
\end{verbatim}
}

\subsubsection*{Evidencia de ejecución}

Salida literal de \texttt{Demo}. Las partes 1 a 6 comprueban el comportamiento diferenciado;
las partes 7 y 8, la restricción de concurrencia.

{\footnotesize
\begin{verbatim}
1. Misma transaccion, cuatro comportamientos (EjecutarTransaccion)
------------------------------------------------------------------
Retiro de 300.00 sobre cada cuenta:
  AHO-001     APLICADA  saldo=     700.00 cargos=  0.00  sin cargos
  COR-001     APLICADA  saldo=    -105.00 cargos=  5.00  comision por sobregiro
  INV-001     APLICADA  saldo=   9,694.00 cargos=  6.00  penalizacion del 2% no reembolsable
  EMP-001     APLICADA  saldo= 199,700.00 cargos=  0.00  registrada
  EMP-002     APLICADA  saldo= 199,700.00 cargos=  0.00  registrada

2. Reglas propias de cada tipo
------------------------------
Retiro de 60,000.00 en empresarial:
  EMP-001     RECHAZADA saldo= 199,700.00 cargos=  0.00  exige doble autorizacion
  EMP-002     APLICADA  saldo= 139,700.00 cargos=  0.00  reportable al regulador
Tres retiros mas en ahorro (limite del mes: 3):
  AHO-001     APLICADA  saldo=     600.00 cargos=  0.00  sin cargos
  AHO-001     APLICADA  saldo=     500.00 cargos=  0.00  sin cargos
  AHO-001     APLICADA  saldo=     398.50 cargos=  1.50  comision por exceso de retiros
Pago de servicios de 15,000.00 en empresarial:
  EMP-002     APLICADA  saldo= 124,697.50 cargos=  2.50  reportable al regulador

3. RevertirTransaccion: tampoco se revierte igual en cada cuenta
----------------------------------------------------------------
  AHO-001     APLICADA  saldo=     698.50 cargos=  0.00  revertida, comision reembolsada
  COR-001     APLICADA  saldo=     200.00 cargos= -5.00  revertida, sobregiro reembolsado
  INV-001     APLICADA  saldo=   9,994.00 cargos=  0.00  penalizacion de 6.00 retenida
  EMP-001     RECHAZADA saldo= 199,700.00 cargos=  0.00  la reversion exige doble firma
  EMP-002     APLICADA  saldo= 124,997.50 cargos=  0.00  revertida con traza de autorizacion

4. AuditarTransaccion: visitante sin estado, informe inmutable
--------------------------------------------------------------
  AHO-001     Ahorro      riesgo=BAJO   asientos=  5  sin hallazgos
  COR-001     Corriente   riesgo=BAJO   asientos=  2  sin hallazgos
  INV-001     Inversion   riesgo=MEDIO  asientos=  2  penalizaciones acumuladas: 6.00
  EMP-001     Empresarial riesgo=BAJO   asientos=  1  sin hallazgos
  EMP-002     Empresarial riesgo=ALTO   asientos=  4  2 operaciones reportables

5. Operacion nueva sin tocar ninguna cuenta (Open/Closed)
---------------------------------------------------------
  AHO-001     Ahorro      interes anual = 3.49
  COR-001     Corriente   interes anual = 0.00
  INV-001     Inversion   interes anual = 399.76
  EMP-001     Empresarial interes anual = 2,396.40
  EMP-002     Empresarial interes anual = 1,499.97

6. Transferencia: dos cerrojos en orden total por identificador
---------------------------------------------------------------
  COR-001     APLICADA  saldo=     150.00 cargos=  0.00  sin cargos
  origen COR-001 -> 150.00 | destino AHO-001 -> 748.50

7. Restriccion: concurrencia sobre la MISMA cuenta
--------------------------------------------------
  con decorador transaccional: esperado 20,000.00, obtenido 20,000.00
  atomico y consistente: true
  sin decorador:               esperado 20,000.00, obtenido 5,440.00
  actualizaciones perdidas:    14,560.00

8. Restriccion: paralelismo sobre cuentas DISTINTAS
---------------------------------------------------
  PAR-AHO     saldo=  10,000.00 correcto=true
  PAR-COR     saldo=  10,000.00 correcto=true
  PAR-INV     saldo=  10,000.00 correcto=true
  PAR-EMP     saldo=  10,000.00 correcto=true
  16 hilos sobre 4 cuentas, sin contencion entre ellas: 10 ms
\end{verbatim}
}

En la parte 7, ocho hilos hacen 250 depósitos de 10,00 sobre la misma cuenta. Con el decorador
el saldo final es exacto; sin él se pierden actualizaciones porque la secuencia leer, sumar y
escribir el saldo deja de ser atómica. La cifra concreta de pérdida cambia en cada ejecución,
porque depende del entrelazado de los hilos: es un error no determinista, que es justamente lo
que lo hace peligroso. Para que la carrera se manifieste de forma visible en cada corrida, la
clase \texttt{Cuenta} amplía a propósito la ventana entre la lectura y la escritura mediante un
\texttt{Thread.yield()} activado solo en esa prueba.

\subsubsection*{Dos observaciones que salieron del experimento}

\textbf{La reversión no es el inverso exacto de la ejecución.} En la cuenta de inversión, el
retiro de 300,00 cobró 6,00 de penalización y la reversión devolvió el principal pero retuvo
la penalización: el saldo no vuelve a 10\,000,00 sino a 9\,994,00. Es una regla de negocio
legítima, y muestra por qué el \textit{rollback} se modela mejor con Command y Memento que
como una resta: el visitante que revierte necesita saber qué ocurrió, no solo qué se pidió. En
la implementación lo consigue buscando el asiento en el historial de la cuenta.

\textbf{Visitor protege frente a comportamiento nuevo, no frente a estado nuevo.}
\texttt{CalcularInteresesVisitor} se añadió sin modificar ninguna cuenta porque calcula a
partir del estado que las cuentas ya exponen. Una operación como ``congelar cuenta''
necesitaría un atributo nuevo en el elemento, y entonces sí habría que tocarlo. Conviene
declararlo: el patrón cierra las clases de los elementos ante operaciones nuevas, no ante
atributos nuevos.

\subsubsection*{Cómo compilar y ejecutar}

{\small
\begin{verbatim}
cd docs/CP4/codigo/ejercicio12
javac -d /tmp/ej12 *.java
java -cp /tmp/ej12 Demo
\end{verbatim}
}

Compila sin advertencias con \texttt{javac -Xlint:all} en OpenJDK 17.

\fussy
